\section{The full correlation range at fixed centre and marginal}
\label{sec:jointrange}
The upper envelope has a complementary sharp lower bound. Together they
determine every possible value of the correlation norm, not merely its
largest value. This also identifies a region in which a nonzero centre
forces a strictly positive measured correlation.

\begin{theorem}[Complete joint range]\label{thm:jointrange}
Fix $0\leq c\leq1$ and $0\leq r<1$. Among all two-qubit density matrices
with Alice steering-centre length $c$ and Bob Bloch length $r$, the set of
attainable squared correlations is exactly
\begin{equation}\label{eq:jointinterval}
\left[\bigl(\max\{0,c(1+r)-1\}\bigr)^2,\;
 (1-r^2)(3-4c)+\bigl(c+(1-c)r\bigr)^2\right].
\end{equation}
Both endpoints are attained. The lower endpoint is attained by a state
diagonal in a product basis. Every lower-endpoint state is locally unitarily
equivalent to the displayed diagonal witness below. At $r=1$, the entire
range is instead the singleton $\{c^2\}$.
\end{theorem}

\begin{proof}
Choose unit vectors $\nn,\mathbf m$ so that the centre is $c\nn$ and
$\bb=r\mathbf m$. When a vector vanishes, choose its direction arbitrarily;
the argument does not divide by $c$ or $r$. Put
$q=\nn^{\mathsf T}T\mathbf m$. The probability of obtaining both minus
outcomes in the corresponding local spin measurements is nonnegative:
\[
 1-\aa\cdot\nn-r+q\geq0.
\]
The centre relation gives
$\aa\cdot\nn=(1-r^2)c+rq$, hence
\begin{equation}\label{eq:lowerfactor}
0\leq(1-r)\bigl(q-[c(1+r)-1]\bigr).
\end{equation}
Since $r<1$, $q\geq c(1+r)-1$. Cauchy--Schwarz gives
$|q|\leq\norm T_F$, proving the lower bound. The upper bound is the sharp
joint envelope proved earlier.

For attainment, put
\[
x=\max\{0,c(1+r)-1\},\qquad a_0=c(1-r^2)+rx.
\]
Take $\aa=a_0\mathbf e_z$, $\bb=r\mathbf e_z$, and
$T=\operatorname{diag}(0,0,x)$. The associated density matrix is
\begin{equation}\label{eq:jointlowerwitness}
\rho^-_{c,r}=\frac14\operatorname{diag}
 (1+a_0+r+x,\;1+a_0-r-x,\;1-a_0+r-x,\;1-a_0-r+x).
\end{equation}
If $x=0$, then $0\leq a_0\leq1-r$, so all entries are nonnegative.
If $x>0$, then $0<x\leq r$, $a_0=1-r+x$, and its eigenvalues are
\[
\frac{1+x}{2},\quad\frac{1-r}{2},\quad\frac{r-x}{2},\quad0.
\]
Its centre is exactly $c\mathbf e_z$ and $F_3^2=x^2$.
These checks include all vanishing diagonal entries, without assuming a
strictly positive state.

Let $\rho^+_{c,r}$ be the upper witness in
Eq.~\eqref{eq:globalwitness}, in the same coordinate convention.
Both witnesses have the same $\bb=r\mathbf e_z$ and the same vector
$\aa-T\bb=(1-r^2)c\mathbf e_z$. Every convex combination
\[
\rho_\alpha=(1-\alpha)\rho^-_{c,r}+\alpha\rho^+_{c,r},
\qquad0\leq\alpha\leq1,
\]
therefore has the same centre and Bob marginal. Write $d=1-c$ and
$z=d+cr$. Direct computation gives
\begin{equation}\label{eq:jointinterpolate}
F_3(\rho_\alpha)^2=
2\alpha^2d(1-r^2)+\bigl[x+\alpha(z-x)\bigr]^2.
\end{equation}
Here $x\geq0$ and $z-x\geq0$; for $x>0$, the latter difference is
$2(1-c)$. Thus the expression is continuous and nondecreasing from the
lower endpoint to the upper endpoint. It supplies every intermediate value.

For the equality classification, if $x>0$ then equality in the lower bound
forces $q=x$ and equality in the Frobenius Cauchy--Schwarz inequality.
Consequently $T=x\nn\mathbf m^{\mathsf T}$, and the centre relation fixes
$\aa=a_0\nn$. If $x=0$, equality means $T=0$, and the centre relation again
fixes $\aa=c(1-r^2)\nn$. Independent local changes of spin axes give precisely
Eq.~\eqref{eq:jointlowerwitness}. Finally, when $r=1$, positivity forces
$\rho=\rho_A\otimes|\beta\rangle\langle\beta|$; its point centre has
length $c$, and its squared correlation is $c^2$.
\end{proof}

The zero-correlation region is exactly $c(1+r)\leq1$.
A state with $T=0$ need not be a product: when both local vectors are
nonzero, the product state would have $T=\aa\bb^{\mathsf T}$.
The lower witness is nevertheless separable because it is diagonal in a
product basis.

For example, at $c=3/4$ and $r=1/2$, the lower witness is
\[
\rho^-_{3/4,1/2}=\operatorname{diag}(9/16,1/4,3/16,0),
\qquad F_3^2=1/64,
\]
while the upper value is $49/64$. All values between them occur at exactly
the same centre length and Bob marginal. At $c=1$, the two endpoints
coincide at $r^2$; the allowed states have pure Alice marginal and are
products. At $r=0$, the range reduces to $[0,(1-c)(3-c)]$.

\begin{corollary}[A quantitative directional consequence]
For every state with $r<1$, centre $c\nn$, and Bob vector $r\mathbf m$,
\[
\nn^{\mathsf T}T\mathbf m\geq c(1+r)-1,
\qquad c(1+r)\leq1+\norm T_F.
\]
If $x=c(1+r)-1>0$ and $F_3^2-x^2\leq\eta$, then
\begin{equation}\label{eq:lowerstability}
\norm{T-x\nn\mathbf m^{\mathsf T}}_F^2\leq\eta,
\qquad
\norm{\aa-a_0\nn}^2\leq r^2\eta.
\end{equation}
\end{corollary}
\begin{proof}
The first assertions were proved above. Since $q\geq x>0$,
\[
\norm{T-x\nn\mathbf m^{\mathsf T}}_F^2
=\norm T_F^2-2xq+x^2\leq\norm T_F^2-x^2\leq\eta.
\]
Subtracting the common centre equations gives
$\aa-a_0\nn=r(T-x\nn\mathbf m^{\mathsf T})\mathbf m$;
the Frobenius bound proves the second assertion.
\end{proof}

This yields an explicit state-distance consequence in the same axes. For
$\rho^-$ with the same centre and Bob vector,
\[
\norm{\rho-\rho^-}_2^2
=\tfrac14\bigl(\norm{\aa-a_0\nn}^2+
\norm{T-x\nn\mathbf m^{\mathsf T}}_F^2\bigr)
\leq\frac{1+r^2}{4}\eta,
\qquad
\norm{\rho-\rho^-}_1^2\leq(1+r^2)\eta.
\]
The norm $\norm{\cdot}_2$ is the Hilbert--Schmidt norm. The same conclusion
holds in the $x=0$ region with $\eta=F_3^2$, since then the correlation
matrix of the witness vanishes. Thus lower-endpoint rigidity has a direct
square-root modulus, including the interface $c(1+r)=1$.
